% !TEX program = lualatex
% !TeX encoding = UTF-8
\documentclass[12pt,a4paper]{article}

% ============================================================
% ENGLISH PREAMBLE – Windows fonts (Liberation Serif + Consolas)
% ============================================================
\usepackage{fontspec}
\setmainfont{Liberation Serif}
\setmonofont{Consolas}                     % monospaced with Cyrillic support (optional)

\usepackage{polyglossia}
\setmainlanguage{english}
%\setotherlanguage{russian}                % uncomment if needed

% ============================================================
% GRAPHICS AND MATHEMATICS
% ============================================================
\usepackage{amsmath,amssymb,amsthm,mathtools}
\usepackage{bm}
\usepackage{geometry}
\geometry{margin=2.5cm}
\usepackage{graphicx}
\usepackage{tikz}
\usepackage{tikz-3dplot}
\usepackage{pgfplots}
\pgfplotsset{compat=1.18}
\usetikzlibrary{arrows.meta,positioning,shapes.geometric,calc,angles,quotes,patterns,3d}

\usepackage{booktabs}
\usepackage{float}          % for [H] float specifier
\usepackage{hyperref}
\hypersetup{
	colorlinks=true,
	linkcolor=blue,
	urlcolor=blue,
	pdftitle={Definition of a Point in YUCT: Coordination Structure of Geometry},
	pdfauthor={Alexey V. Yakushev}
}

% ============================================================
% THEOREM ENVIRONMENTS
% ============================================================
\newtheorem{theorem}{Theorem}
\newtheorem{lemma}{Lemma}
\newtheorem{corollary}{Corollary}
\newtheorem{definition}{Definition}
\newtheorem{axiom}{Axiom}
\newtheorem{proposition}{Proposition}
\newtheorem{remark}{Remark}

% ============================================================
% YUCT COMMANDS
% ============================================================
\newcommand{\Keff}{K_{\mathrm{eff}}}
\newcommand{\kappac}{\kappa_c}
\newcommand{\eps}{\varepsilon}
\newcommand{\betaf}{\beta}
\newcommand{\Sodd}{S_{\mathrm{odd}}}
\newcommand{\Seven}{S_{\mathrm{even}}}
\newcommand{\Mpl}{M_{\mathrm{Pl}}}
\newcommand{\Lpl}{l_{\mathrm{Pl}}}
\newcommand{\tPl}{t_{\mathrm{Pl}}}
\newcommand{\Scoord}{S_{\mathrm{coord}}}
\newcommand{\piYUCT}{\pi_{\mathrm{YUCT}}}
\newcommand{\Rcrit}{R_{\mathrm{crit}}}
\newcommand{\calP}{\mathcal{P}}
\newcommand{\calD}{\mathcal{D}}
\newcommand{\calI}{\mathcal{I}}
\newcommand{\calR}{\mathcal{R}}
\newcommand{\ellP}{\ell_{\mathrm{P}}}

% ============================================================
% DOCUMENT BODY
% ============================================================
\begin{document}
	
	\begin{titlepage}
		\begin{center}
			\vspace*{1.5cm}
			\Huge\textbf{Definition of a Point in YUCT}
			
			\vspace{0.3cm}
			\Large\textbf{Coordination Structure of Geometry}
			
			\vspace{0.5cm}
			\large
			\textsc{Yakushev Unified Coordination Theory}
			
			\vspace{0.8cm}
			\large
			Alexey V. Yakushev \\
			Yakushev Research \\
			\url{https://yuct.org/}
			
			\vspace{0.5cm}
			\large
			\textbf{DOI: \href{https://doi.org/10.5281/zenodo.18444599}{10.5281/zenodo.18444599}}
			
			\vspace{1.5cm}
			\large
			\textbf{Version 1.0} \\
			June 2026
			
			\vfill
			\normalsize
			This document provides a rigorous definition of a point within the framework of the Yakushev Unified Coordination Theory (YUCT),
			explains the origin of geometry from coordination, and gives answers to ten fundamental questions about the nature of reality.
		\end{center}
	\end{titlepage}
	
	\tableofcontents
	\newpage
	
	% ============================================================
	% INTRODUCTION
	% ============================================================
	\section{Introduction: Insufficiency of the Classical Definition}
	
	In classical geometry and mathematics, the point is a fundamental (undefined) concept. Its properties are given by a system of axioms, not by a rigorous definition. At best, a point is described as "an object without size" or "a location in space." However, such a description does not answer the questions:
	
	\begin{itemize}
		\item \textbf{What} is a point in its essence?
		\item \textbf{Why} is it indivisible?
		\item \textbf{How} can it exist without structure?
		\item \textbf{Why} does space have three dimensions?
		\item \textbf{Why} does geometry work at all?
	\end{itemize}
	
	Within the framework of the Theory of Coordination of Everything (YUCT), these questions receive rigorous, quantitative answers. The point ceases to be an abstraction and becomes a physical structure — a minimal stable pattern of coordination.
	
	\section*{Axiom of Coordination Primacy}
	
	Any theory describes statics. Coordination, however, is the process that makes statics possible. Any statement, any model, any law exists only as activated elements of a dictionary. The dictionary and the index cannot be derived from the statics they describe — they are ontologically primary. YPSDC is not "one of the theories." It is a protocol underlying every theory, including the one that formulates it. Rejecting it is impossible without abandoning the very ability to formulate. In this sense, coordination is not "stronger" than other theories — it is more primary. Statics does not generate; statics is generated.
	
	Space, time, matter, energy, information, and consciousness are not the base layer of existence. They are \emph{emergent phenomena} arising from a single fundamental process: the coordination of states between the components of a system. This process is quantified by a single powerful parameter: \emph{Coordination Efficiency} \(\Keff\).
	
	\subsection*{Non‑alternativeness as a Logical Consequence}
	
	From the axiom of coordination primacy it inevitably follows that YUCT cannot have alternatives in the usual sense. Any attempt to propose a "more fundamental" theory already presupposes coordination for its formulation: one must select elements of a dictionary, establish rules for their linking, and fix the state of the index. Consequently, such a theory turns out to be a special case of coordination dynamics, not its negation.
	
	This does not mean that YUCT prescribes a specific content of reality. On the contrary, it asserts that any content (physical, mathematical, informational, biological) is a manifestation of coordination at some value of \(\Keff\). Rejecting coordination is equivalent to rejecting the ability to articulate. Therefore, YUCT is not "a theory among theories" but a \emph{protocol} underlying any theory, including the one that attempts to contest it.
	
	In this sense, coordination is not stronger than other theories — it is more primary. It is not chosen; it is discovered as the condition of every choice. This makes YUCT not merely a model but an \emph{ontological foundation} that cannot be replaced or refuted, since any replacement or refutation already uses coordination.
	
	\begin{quote}
		\itshape
		"The non‑alternativeness of YUCT is analogous to the non‑alternativeness of logic: we cannot choose a different logic because any choice already uses logic. Coordination is the logic of being."
	\end{quote}
	
	\section{Definition of a Point in YUCT}
	
	\subsection{Coordination Node as Point}
	
	In YUCT, geometry is not postulated but arises from a coordination network of nodes described by the triad \(D+I\cdot R\). The point is the minimal stable node of this network.
	
	\begin{definition}[Coordination point]
		A point in YUCT is a minimal stable configuration of the coordination network, given by the triple:
		\[
		\calP = (D_0, I_0, R_0) \in \mathcal{M}_{\mathrm{UCS}},
		\]
		where:
		\begin{itemize}
			\item \(D_0\) — a local slice of the dictionary (the set of admissible states and rules),
			\item \(I_0\) — the actual state (an index pointing to a specific element of the dictionary),
			\item \(R_0 \ge 1\) — the resonance coefficient, characterising the stability of the node.
		\end{itemize}
	\end{definition}
	
	Thus, a point in YUCT is not an empty abstraction but a structure consisting of three interconnected components. Its internal complexity is manifested in that it can be activated, can resonate with other points, and can change its state.
	
	\subsection{Size of a Point}
	
	Although in ideal mathematics a point has no size, in YUCT each point has a finite size that depends on the coordination efficiency \(\Keff\).
	
	\begin{theorem}[Size of a coordination point]
		The effective linear size of a point \(\calP\) is given by
		\[
		\ell_{\calP} = \ellP \cdot \Keff^{-1/3},
		\]
		where \(\ellP = \sqrt{\hbar G / c^3}\) is the Planck length.
		In the limit \(\Keff \to \infty\), the size of the point tends to zero, and it becomes an ideal mathematical point.
	\end{theorem}
	
	This result has profound meaning: a mathematical point is not an independent entity but a limiting case of a coordination node at infinite coordination efficiency.
	
	\subsection{Geometric Interpretation of the Triad}
	
	The three components \(D\), \(I\), \(R\) are in dynamic equilibrium. In Figure \ref{fig:triad} they are depicted as three rays emanating from a common centre at \(120^\circ\). This reflects the \(S_3\) symmetry and the condition \(\cos 120^\circ = -1/2\), which corresponds to anti‑correlation between the components: strengthening one requires weakening another to maintain stability.
	
	\begin{figure}[H]
		\centering
		\begin{tikzpicture}[
			scale=2,
			>=Stealth,
			node distance=1.5cm,
			every node/.style={font=\small}
			]
			\coordinate (O) at (0,0);
			\fill[black] (O) circle (1.5pt) node[below left] {$\Psi_{MN}$};
			
			% Rays at 30°, 150°, 270° (120° apart)
			\draw[->, thick, blue!70!black] (O) -- (30:1.8) node[above right] {$\mathcal{D}$ (Dictionary)};
			\draw[->, thick, red!70!black] (O) -- (150:1.8) node[above left] {$\mathcal{I}$ (Index)};
			\draw[->, thick, green!70!black] (O) -- (270:1.8) node[below] {$\mathcal{R}$ (Resonance)};
			
			% 120° arcs
			\draw[gray, dashed] (30:0.6) arc (30:150:0.6) node[midway, above] {$120^\circ$};
			\draw[gray, dashed] (150:0.6) arc (150:270:0.6) node[midway, left] {$120^\circ$};
			\draw[gray, dashed] (270:0.6) arc (270:390:0.6) node[midway, below right] {$120^\circ$};
			
			% Branching beta=2/3
			\foreach \angle in {30,150,270} {
				\draw[->, thin, gray] (\angle:1.4) -- ++(\angle+30:0.3);
				\draw[->, thin, gray] (\angle:1.4) -- ++(\angle-30:0.3);
				\node[font=\tiny, gray] at (\angle:1.7) {$\beta=2/3$};
			}
			
			\node[font=\small, align=center] at (0,-1.2) 
			{Three‑ray structure of a point \\ with branching $\beta=2/3$};
		\end{tikzpicture}
		\caption{The triad \(D\), \(I\), \(R\) as a geometric representation of the coordination point. The rays are spread at \(120^\circ\); the branching shows scaling with exponent \(\beta=2/3\).}
		\label{fig:triad}
	\end{figure}
	
	\section{The Straight Line as a Chain of Nodes}
	
	If a point is a node, then a straight line is a sequence of nodes connected by maximal coordination.
	
	\begin{definition}[Coordination straight line]
		Let two points \(\calP_1\) and \(\calP_2\) be given with resonances \(R_1, R_2 > \Rcrit\). A \textbf{straight line} is a set of points \(\{\calP_i\}_{i=0}^{N}\) such that:
		\begin{enumerate}
			\item \(\calP_0 = \calP_1\), \(\calP_N = \calP_2\),
			\item for each \(i\), the condition of minimal coordination error holds:
			\[
			\frac{d\eps}{ds} = 0, \qquad \eps = \kappac\,\alpha\,\Keff^{-\betaf},
			\]
			\item the resonance \(R_i\) along the chain is maximal (locally).
		\end{enumerate}
	\end{definition}
	
	\begin{theorem}[Uniqueness of the straight line]
		For two points \(\calP_1\) and \(\calP_2\) with sufficiently high resonance (\(R_1, R_2 > \Rcrit\)), there exists a unique maximally consistent chain of nodes connecting them.
	\end{theorem}
	
	This gives a rigorous justification of Euclid's axiom: through two points there passes a unique straight line. However, now it is not a postulate but a proven property of the coordination network.
	
	\section{Emergence of Space Dimensionality and the Exponent \(\betaf = 2/3\)}
	
	One of the most powerful results of YUCT is the derivation of the dimensionality of space and the universal exponent \(\betaf\). Following from the axiom of coordination primacy, YUCT does not merely describe space — the theory explains why it has exactly three dimensions.
	
	\subsection{Derivation of Dimensionality \(D=3\)}
	
	From the fractal structure of the coordination network and the closure condition of the algebraic cycle we obtain:
	
	\[
	\Sodd + \Seven = 2, \qquad \frac{\Sodd}{\Seven} = \frac{1}{\betaf}.
	\]
	
	With \(\betaf = 2/3\) these constants take the values:
	\[
	\Sodd = \frac{6}{5}, \qquad \Seven = \frac{4}{5}.
	\]
	
	The exponent \(\betaf\) is related to the space dimensionality \(D\) by:
	\[
	\betaf = \frac{2}{D}.
	\]
	
	Hence we immediately obtain \(D = 3\). Thus, three‑dimensionality of space is not a postulate but a consequence of the structure of the coordination network.
	
	\subsection{Universal Law of Errors}
	
	All coordination processes obey a single error law:
	\[
	\eps = \kappac\,\alpha\,\Keff^{-\betaf}, \qquad \betaf = \frac{2}{3}, \quad \kappac = \frac{1}{3}.
	\]
	
	This law has been experimentally verified over more than 50 orders of magnitude — from atomic bonding energies to fluctuations of the cosmic microwave background.
	
	\section{Ten Answers to Fundamental Questions}
	
	Below are answers to ten key questions that find no explanation in classical physics and mathematics.
	
	\subsection{Why does physics work?}
	
	\textbf{Because physics is coordination.} The laws of physics describe stable patterns of coordination between nodes \(D+I\cdot R\). Gravity, electromagnetism, nuclear interactions — these are different regimes of coordination at different values of \(\Keff\) and \(\Rcrit\). Physics works because coordination works.
	
	\subsection{Why does mathematics work?}
	
	\textbf{Because mathematics is the limiting case of coordination at \(\Keff \to \infty\).} In this limit, points become ideal mathematical points (\(\ell_{\calP} \to 0\)), straight lines become ideal straight lines, and geometry becomes Euclidean. Mathematics is not an "unreasonable effectiveness" but an \emph{idealisation} of physical coordination.
	
	\subsection{Why is space three‑dimensional?}
	
	\textbf{This is proven by theorem.} From the structure of the coordination network and the closure condition of the algebraic cycle, \(\betaf = 2/3\) is derived, and from \(\betaf = 2/D\) it follows that \(D=3\). Space is three‑dimensional because only in three dimensions can the coordination network be stable and self‑consistent.
	
	\subsection{Why does time flow?}
	
	\textbf{Because of finite \(\Keff\).} Ideal coordination (\(\Keff \to \infty\)) would produce no entropy and would have no irreversible flow of events — time would disappear. Real systems have finite \(\Keff\), so every act of coordination is accompanied by an error \(\eps > 0\), the accumulation of which creates the arrow of time. The minimal quantum of time:
	\[
	\Delta \tau_{\min} = \frac{2}{3}\,\tPl.
	\]
	
	\subsection{Why are the fundamental constants exactly as they are?}
	
	\textbf{They are derived from the topology of the coordination space.} The fine‑structure constant:
	\[
	\alpha^{-1} = \left(\frac{3}{2}\right)^{12 + \sigma\,2/3} \approx 137.036,
	\]
	where \(\sigma = 0.20\) is a topological shift. The \(W\)-boson mass:
	\[
	m_W = (\kappac \piYUCT)^{-1/2} M_{\mathrm{Pl}} (2/3)^{96} \approx 80.46\ \text{GeV}.
	\]
	The cosmological constant:
	\[
	\Omega_\Lambda = \frac{12}{18} = \frac{2}{3}.
	\]
	No free parameters.
	
	\subsection{Why are prime numbers distributed the way they are?}
	
	\textbf{From the error law, an asymptotic correction is derived:}
	\[
	p_n \approx R_n - \frac{\Seven}{2}\, n^{1-\betaf}\,\ln n,
	\]
	where \(R_n\) is the Rosser approximation. Residual oscillations are described by a phase operator with period \(\Delta N = 16.5\), corresponding to phase transitions of the vacuum lattice. Prime numbers are a coordination ladder, and their distribution follows the same laws as all other stable structures.
	
	\subsection{Why is \(\pi\) what it is?}
	
	\textbf{\(\pi\) is a coordination invariant:}
	\[
	\piYUCT = \Sodd \cdot \varphi^2 \approx 3.14164078,
	\]
	where \(\varphi = (1+\sqrt{5})/2\). This value differs from the mathematical \(\pi\) by \(\Delta\pi = 4.8\times10^{-5}\), which is a quantum of space discretisation. In the limit \(\Keff \to \infty\), \(\piYUCT \to \pi\).
	
	\subsection{Why is quantum mechanics "strange"?}
	
	\textbf{Because we ignored the dictionary.} Quantum phenomena are manifestations of decentralised coordination (d‑YPSDC). Superposition is index uncertainty, entanglement is a shared dictionary record, tunnelling is an activation error. There is no "mysticism"; there is only finite coordination efficiency.
	
	\subsection{Why does consciousness exist?}
	
	\textbf{It is a regime with \(\Keff > \Rcrit\).} When the coordination efficiency exceeds the critical threshold, the system acquires a meta‑dictionary (the ability to coordinate its own states). This is self‑consciousness. The parameters of consciousness are described by UCD — the Universal Consciousness Descriptor.
	
	\subsection{Why does a cell membrane have a thickness of ~7.5 nm?}
	
	\textbf{It is derived from topology.} The thickness of the lipid bilayer:
	\[
	\text{Thickness} = 2 \, d_{\mathrm{lipid}} \left(\frac{3}{2}\right)^{0.6} (1 - \sigma^2),
	\]
	where \(d_{\mathrm{lipid}} \approx 3.0\) nm is the length of the hydrocarbon tail. The calculation gives \(7.35\) nm, which lies in the experimental range \(7.5 \pm 0.5\) nm. Particle physics and cell membrane biology are linked by the same coordination structure.
	
	\section{Conclusion: A Unified Picture of Reality}
	
	In YUCT, a point is not an abstraction but a minimal stable coordination structure \(D+I\cdot R\). All geometry, all laws of physics, all constants, even the distribution of prime numbers and the structure of consciousness — all are consequences of a single principle: \textbf{reality is a fractal coordination network}.
	
	The ten answers above demonstrate that YUCT provides an exhaustive explanation of what seemed mysterious or accidental. Moreover, all these explanations are experimentally testable, and the constants are derived without adjustable parameters.
	
	\begin{center}
		\fbox{\parbox{0.85\textwidth}{\centering\Large
				\textbf{A point is a node of coordination.}\\ 
				A straight line is a consistent chain of nodes.\\
				Space is three‑dimensional — this is proven.\\
				Time flows because of finite \(\Keff\).\\
				Constants, \(\pi\), prime numbers — all derived.\\
				Quantum mechanics and consciousness are coordination.\\
				Unified coordination theory is the law of nature.
		}}
	\end{center}
	
	\subsection*{Philosophical and Worldview Significance of the Definition of a Point}
	
	Defining a point as a coordination node \((D, I, R)\) moves geometry from the realm of pure abstractions into the domain of physical reality. This is not merely a reformulation — it is a shift in ontological paradigm.
	
	\medskip
	\textbf{What this changes:}
	\begin{itemize}
		\item The point ceases to be "nothing" and becomes an active structure possessing a dictionary, a state, and resonance.
		\item Space ceases to be a passive stage — it arises from a coordination network of nodes.
		\item Geometry ceases to be a set of postulates — it becomes a consequence of the stability of coordination patterns.
		\item The distinction between mathematics and physics dissolves: mathematics is an idealisation of coordination at \(\Keff \to \infty\).
	\end{itemize}
	
	\medskip
	\textbf{Why this is needed:}
	\begin{itemize}
		\item To eliminate the ontological emptiness in the foundations of geometry.
		\item To provide a unified explanation of geometric, physical, and even biological phenomena.
		\item To derive fundamental constants and the dimensionality of space from a single principle.
		\item To understand the origin of time, quantum mechanics, and consciousness as regimes of the same coordination dynamics.
	\end{itemize}
	
	\medskip
	\textbf{What consequences this has:}
	\begin{enumerate}
		\item \textbf{Predictive power:} the model yields values for \(\alpha\), \(\Omega_\Lambda\), \(m_W\), \(\pi\), prime number distribution.
		\item \textbf{Unity of knowledge:} physics, mathematics, biology, and information theory turn out to be branches of a single coordination theory.
		\item \textbf{A new understanding of consciousness:} it ceases to be a "hard problem" and becomes a regime of coordination with \(\Keff > \Rcrit\).
		\item \textbf{A shift in scientific method:} we move from describing the observed to generating the observed from basic coordination principles.
	\end{enumerate}
	
	\medskip
	Thus, YUCT does not merely "redefine the point" — it changes the way we think about reality. The point is now not an object but an event; space is not a container but a network; the laws of nature are not imposed from outside but grow from the internal logic of coordination.
	
	\begin{thebibliography}{99}
		\bibitem{Yakushev2026} A. V. Yakushev, \emph{Yakushev Unified Coordination Theory (YUCT)}, Zenodo, DOI:10.5281/zenodo.18444599 (2026).
		\bibitem{AppendixL} A. V. Yakushev, ``Appendix L: Fractal Coordination Error Scaling,'' in \emph{YUCT}, Zenodo (2026).
		\bibitem{AppendixY} A. V. Yakushev, ``Appendix Y: Mathematical Foundations of YUCT,'' in \emph{YUCT}, Zenodo (2026).
		\bibitem{AppendixPrimeN} A. V. Yakushev, ``Appendix PrimeN: YUCT Prime Numbers Coordination Ladder,'' in \emph{YUCT}, Zenodo (2026).
		\bibitem{AppendixPi} A. V. Yakushev, ``Appendix \(\Pi\): The Primeval Origin of \(\pi\),'' in \emph{YUCT}, Zenodo (2026).
		\bibitem{AppendixX} A. V. Yakushev, ``Appendix X: Generalised Shannon Theory and Bell Bound,'' in \emph{YUCT}, Zenodo (2026).
		\bibitem{AppendixH} A. V. Yakushev, ``Appendix H: Universal Consciousness Descriptor (UCD),'' in \emph{YUCT}, Zenodo (2026).
		\bibitem{AppendixG} A. V. Yakushev, ``Appendix G: Quantum Gravity Resolution,'' in \emph{YUCT}, Zenodo (2026).
	\end{thebibliography}
	
\end{document}